\section{Introduction}

The severe acute respiratory syndrome coronavirus-2 (SARS-CoV-2) and the associated COVID-19 pandemic have created a substantial and unforeseen burden on the global healthcare system \cite{abrams2020ecmo}. With a global mortality of over 6.8 million (as of Feb 2023), there is considerable focus on therapeutic solutions for patients with most severe manifestations of the disease. In many cases, scarce treatment options need to be considered and evaluated to support patients' lives. In particular, World Health Organization recommends extracorporeal membrane oxygenation (ECMO) for patients who are refractory to conventional therapies, a support modality only available in expert centres with sufficient experience \cite{world2021covid}. As such treatments are technically complex and resource-intensive with difficulty in predicting outcomes, proper treatment evaluation and assignment for ECMO has been the subject of significant debate since the start of the pandemic \cite{falcoz2020extracorporeal,haiduc2020role}. 
Recent reports point to the vitalness of ECMO support over 14,000 reported COVID-19 patients with an overall hospital mortality of approximately 47\% \cite{ELSO}. In comparison, nearly 90\% who couldn't find a spot at an ECMO center died, and these patients were young and previously healthy, with a median age of 40 \cite{gannon2022association}. In addition to the vitalness, demand for ECMO far exceeded its availability leading to numerous patients waiting for ECMO support. To date, patient triage and ECMO resource allocation have been limited to the use of Intensive Care Unit (ICU) illness markers and markers of severe medically refractory respiratory failure, neither of which has been validated to predict patients who would ultimately benefit from this resource-intensive, high-risk therapy \cite{salyer2021first,shekar2020extracorporeal}. 
These gaps in knowledge highlight the need to develop clinically applicable predictive models to assist clinicians in identifying patients most likely to benefit from ECMO support and evaluating the treatment effect thus aiding in patient triage and the necessary resource allocation \cite{salyer2021first,huang2021role}.

From the perspective of treatment effect analysis, treatment assignment indicates whether a patient received the ECMO treatment, “factual outcomes” corresponds to patient’s discharge status (i.e. whether the patient survived/died), and patient’s features (or coviatates) are the electronic health records (EHR) before ECMO initiation. To support ECMO treatment decisions, we need to make two-side estimations. First, we need to model the probability of getting treatment for each patient (propensity scoring), to reflect the underlying treatment assignment policy \cite{rubin} as well as the associated risk consideration, as ECMO itself can lead patients to death. Second, we need to estimate the impact of each treatment decision, calculated by the survival/death difference with and without treatment.


Developing an ECMO decision-assistant model differs from typical supervised machine learning problems  or standard treatment effect problems in healthcare. Compared with a supervised clinical problem \cite{xue2021use,abraham2023integrating,jiao2022continuous,li2022self,shi2018obstructive,xue2022distance,xue2022validation,liu2022predicting}, the entire vector of treatment effects can never be obtained, but only the factual outcomes aligned with the individualized treatment assignments. Compared with a typical treatment effect problem \cite{spirtes2009tutorial,rubin}, it faces the challenges of strong selection bias, scarcity of treatment cases and curse of dimensionality. 
First, unlike randomized controlled trials or many common datasets, ECMO prediction is prone to strong selection bias. ECMO is only applied to high-risk patients with severe symptoms, when few life-supporting alternatives are available. As a result, the features characterizing severity are significantly different from those among non-ECMO-treated patients (referred hereafter as control patients). For treatment effect models that apply a supervised learning framework for each treatment option separately, the learned models would not generalize well to the entire population. Second, the technical complexity and resource intensiveness of such treatment limit the number of ECMO assignments, resulting in a much smaller cohort size when compared to control patients.  In fact, a retrospective study shows that fewer than 0.7\% of critically-illed COVID-19 patients received ECMO treatment\cite{bertini2021ecmo}. Moreover, the EHR dataset contains hundreds types of measurements/lab tests, and only a small subset of these measurements/tests would be conducted to each patient. Due to the limited cohort size in ECMO patients, it is challenging for most machine learning models to overcome the curse of dimensionality without falling into over-fitting. Instead of directly capturing the relationship between the prediction tasks and these partially-observed high-dimensional input features, a lower-dimensional representation of inputs is desired \cite{xuebingkdd}.  

In this paper, we tackle these challenges and propose \textit{Treatment Variational AutoEncoder (TVAE)}, a novel approach that uses a disentangled and balanced latent representation to infer a subject’s potential (factual and counterfactual) outcomes and treatment assignment. It leverages the recent advances in representation learning, and extends the capability of deep generative models in the following aspects:
\begin{itemize}
\item \textbf{Treatment Joint Inference}: The lower-dimensional latent representation of TVAE is \textit{semi-supervised} by treatment assignment and factual response. This architecture eliminates the need of auxiliary networks for prediction, and regulates the counterfactual prediction.  
\item \textbf{Distribution Balancing}: To generate an accurate latent representation, the selection bias is delicately disentangled from other latent dimensions to facilitate treatment assignment prediction and maximize the information sharing between two groups. 
\item \textbf{Label Balancing}: Utilizing the generative function of TVAE, we create \textit{fake} ECMO cases from the posterior distribution of \textit{real} ECMO cases, hence addressing the data imbalance and over-fitting without perturbing the patient distribution. 
\end{itemize}

  Our proposed TVAE outperforms state-of-the-art treatment effect models in predicting ECMO treatment assignment as well as factual responses (with or without ECMO), validated by two large real-world COVID-19 datasets consisting of an international dataset including 118,801 Intensive Care Units (ICU) patients from 1651 hospitals and a institutional dataset including 6,016 ICU patients from 15 hospitals. It also achieves the best performance in estimating individual treatment effects using the public IHDP dataset. While this work is motivated by and evaluated in the context of ECMO treatment, the proposed approach may be generalized for other treatment estimation tasks facing selection bias, label imbalance, and curse of dimensionality.


